% =============================================================================
\section{Modelo físico y sistema simulado}
\label{sec:modelo}
% =============================================================================

\subsection{Potencial de Lennard-Jones}

La dinámica del sistema se describe mediante el potencial de Lennard-Jones (LJ)
por pares:
%
\begin{equation}
    V_{\mathrm{LJ}}(r) = 4\varepsilon
    \left[
        \left(\frac{\sigma}{r}\right)^{12}
        -
        \left(\frac{\sigma}{r}\right)^{6}
    \right],
    \label{eq:lj}
\end{equation}
%
donde $r$ es la distancia interatómica, $\eps$ es la profundidad del pozo de
potencial y $\sigma$ es el diámetro efectivo. El término $r^{-12}$ modela la
repulsión de corto alcance (exclusión de Pauli); el término $r^{-6}$ corresponde
a la dispersión de London. Para la interacción proyectil-superficie se introduce
el parámetro de adhesión $\epsps$:
%
\begin{equation}
    V_{12}(r) = 4\epsps
    \left[
        \left(\frac{\sigma}{r}\right)^{12}
        -
        \left(\frac{\sigma}{r}\right)^{6}
    \right].
    \label{eq:lj_ps}
\end{equation}
%
Se exploram $\epsps/\eps = 1{,}0$ (homoatómico) y $0{,}3$ (adhesión débil).
El potencial se trunca en $r_c = 2{,}5\,\sigma$ con corrección de cola de
energía. En \textbf{unidades reducidas de LJ} ($\eps, \sigma, \tau$ como
escalas de energía, longitud y tiempo):
%
\begin{equation}
    T^* = \frac{k_B T}{\eps},
    \qquad
    v^* = \frac{v}{\sigma/\tau},
    \qquad
    E^* = \frac{E}{\eps},
    \label{eq:unidades}
\end{equation}
%
con $\tau = \sigma\sqrt{m/\eps}$. Para argón: $\eps = 1{,}04\times10^{-2}$~eV,
$\sigma = 3{,}4$~\AA{}, $\tau \approx 2{,}15$~ps.

\subsection{Superficie FCC y proyectil}

La superficie es un \textit{slab} de argón FCC de $8\times8\times6$ celdas
unitarias ($\mathcal{N}_{\mathrm{slab}} = 1\,536$ átomos, parámetro de red
$a_0 = 2^{1/6}\sqrt{2}\,\sigma \approx 1{,}5496\,\sigma$), dividido en tres
regiones (Tabla~\ref{tab:regiones}). El proyectil es un clúster icosaédrico de
Mackay de $N$ átomos ($N \in \{1, 2, 7, 13, 55\}$), pre-minimizado a $T^*=0$
mediante minimización de gradiente conjugado seguida de FIRE~\cite{Bitzek2006}.

\begin{table}[ht]
    \centering
    \small
    \caption{Estratificación del \textit{slab} superficial.}
    \label{tab:regiones}
    \begin{tabular}{lll}
        \toprule
        Región & Tratamiento & Función \\
        \midrule
        Inferior ($\approx 2$ celdas) & \texttt{setforce 0} & Simula \textit{bulk} semi-infinito \\
        Intermedia & NVT Nosé-Hoover & Baño de calor a $\Tstar$ \\
        Superior ($1$ celda) & NVE & Superficie libre para impacto \\
        \bottomrule
    \end{tabular}
\end{table}

Esta estratificación impide reflexiones artificiales de las ondas de fonones en
los bordes $z$. Adicionalmente, se aplica una pared reflectante
(\texttt{wall/reflect}) en el techo de la caja para evitar la evaporación de
átomos a temperaturas $\Tstar \geq 0{,}3$.

\subsection{Observables del impacto}
\label{sec:observables}

El \textbf{coeficiente de restitución} (COR) normal es el observable primario:
%
\begin{equation}
    e = \frac{|v_{f,z}|}{|v_{0,z}|}, \qquad 0 \leq e \leq 1,
    \label{eq:cor}
\end{equation}
%
donde $v_{0,z} = -v_0\cos\theta$ y $v_{f,z}$ son las velocidades del CM del
proyectil antes y después del impacto. El caso $e=0$ define el umbral de
\textit{sticking}.

Para $N>1$, la energía cinética del clúster se descompone en:
%
\begin{equation}
    E_k^{\mathrm{total}} = \kecm + \keint,
    \qquad
    \keint = \sum_{i=1}^N \tfrac{1}{2}m|\mathbf{v}_i - \vcm|^2,
    \label{eq:ek_descomp}
\end{equation}
%
donde $\keint$ incluye tanto los modos vibracionales como rotacionales del
clúster. La \textbf{energía incidente normal} es:
%
\begin{equation}
    \Einc = \tfrac{1}{2}N m v_0^2\cos^2\theta.
    \label{eq:einc}
\end{equation}

La \textbf{partición energética} durante el impacto obedece el balance:
%
\begin{equation}
    1 = \cor^2 + \alphaint + \alphasurf,
    \label{eq:balance}
\end{equation}
%
donde $\alphaint = \Delta\keint/\Einc$ es la fracción de energía depositada en
modos internos del clúster, y $\alphasurf = 1 - \cor^2 - \alphaint$ es la
fracción transferida a los fonones de la red (estimada por balance, dado que el
termostato NVT impide una medición directa fiable mediante $\Delta T$).

El \textbf{número de Weber} es la razón entre energía cinética incidente y
energía de adhesión por átomo:
%
\begin{equation}
    \We = \frac{\Einc}{N\,\epsps}
         = \frac{v_0^2\cos^2\theta}{2\,\epsps}.
    \label{eq:weber}
\end{equation}
%
Para $\We \gg 1$ domina el rebote; para $\We \lesssim 1$ el proyectil queda
atrapado (\textit{sticking}). Este número adimensional permite comparar
resultados obtenidos con distintos potenciales y materiales~\cite{Dahneke1975,
Hou2000}.

\subsection{Protocolo de simulación}

\paragraph{Equilibración.} La superficie se equilibra bajo NVT con calentamiento
gradual de $\Tstar = 0{,}01$ a $\Tstar_{\mathrm{target}}$ en 50\,000 pasos,
seguido de 35\,000 pasos de producción. El estado final se guarda en un archivo
\textit{restart} de \lammps{} reutilizado por todas las réplicas.

\paragraph{Impacto.} Se carga el \textit{restart}, se añade el clúster con el
comando \texttt{read\_data\ldots add append} (desplazado a $z_0 = z_{\mathrm{sup}}
+ \Delta z_{\mathrm{gap}}$) y se lanzan 90\,000 pasos de integración con
$\Delta t = 0{,}002\,\tau$. Se registra la cinemática del CM del clúster
($z_{\mathrm{CM}}, v_{z,\mathrm{CM}}, v_{x,\mathrm{CM}}, \keint$) cada 10 pasos.

\paragraph{Espacio de parámetros.} El barrido cubre
$N \in \{1, 2, 7, 13, 55\}$,
$v_0 \in \{0{,}5,\, 1{,}0,\, 2{,}0,\, 3{,}0,\, 5{,}0\}\,\sigma/\tau$,
$\theta \in \{0°, 45°, 60°\}$,
$\Tstar \in \{0{,}1,\, 0{,}3,\, 0{,}5\}$,
$\epsps \in \{0{,}30,\, 1{,}00\}$,
con 3 réplicas por condición, totalizando 1\,350 simulaciones independientes.
